\documentclass[12pt,a4paper,oneside]{article}
\usepackage[brazil]{babel}
\usepackage[utf8]{inputenc}
\usepackage[dvipsnames]{xcolor}
\usepackage{listings}
\usepackage{wrapfig}
\usepackage{olymp}

\setcounter{problem}{1}

\begin{document}

\begin{problem}{Divisores}{divisores}{2}

\begingroup
\setlength{\intextsep}{0pt}%
%\setlength{\columnsep}{0pt}

O seguinte programa foi escrito para informar quantos divisores cada número no intervalo possui.

\lstset{language=C++,
                basicstyle=\ttfamily,
                keywordstyle=\color{Mulberry}\ttfamily,
                stringstyle=\color{Maroon}\ttfamily,
                commentstyle=\color{YellowGreen}\ttfamily,
                morecomment=[l][\color{magenta}]{\#},
                basewidth = {.48em},
                showstringspaces=false
}
{%\small
\begin{lstlisting}
#include <iostream>
using namespace std;

int main() {
  int a, b;
  cin >> a >> b;

  for (int i=a ; i<=b ; i++)
    cout << i << ": " << divisores(i) << endl;

  return 0;
}
\end{lstlisting}
}

O problema é que não existe a função divisores usada neste programa. Sua tarefa é completar o programa criando esta função. A função deve seguir o cabeçalho abaixo.
Ela recebe um número inteiro e retorna quantos divisores ele tem.

\texttt{\textcolor{Mulberry}{int} divisores(\textcolor{Mulberry}{int} n)}

% \Notes

\InputFile

A entrada contém dois valores inteiros, $A$ e $B$.

Restrição: $1 \leq A \leq B \leq 1000000$. O intervalo terá no máximo $20$ valores.

\endgroup

\OutputFile

O programa escreve uma linha na saída para cada valor no intervalo de $A$ a $B$, inclusive eles, informando quantos divisores cada valor tem, conforme mostrado nos exemplos.

% \Notes

\Example

\begin{example}
\exmp{
1 10
}{
1: 1
2: 2
3: 2
4: 3
5: 2
6: 4
7: 2
8: 4
9: 3
10: 4
}%
\end{example}

\begin{example}
\exmp{
20 30
}{
20: 6
21: 4
22: 4
23: 2
24: 8
25: 3
26: 4
27: 4
28: 6
29: 2
30: 8
}%
\end{example}

\begin{example}
\exmp{
1000 1005
}{
1000: 16
1001: 8
1002: 8
1003: 4
1004: 6
1005: 8
}%
\end{example}

%Comentários sobre o exemplo, se necessário.

\end{problem}

\end{document}